%! Solving Rational Equations — Unit 4, Lesson 4.6 — Guided Notes (Answer Key)
\documentclass[10pt]{article}
\usepackage{algebra2-article}
\usepackage{algebra2-key}   % pulls in -boxes; do NOT also load -boxes
\newcommand{\TallMath}[1]{$\displaystyle #1\rule[-1.4em]{0pt}{3.2em}$}
% Standard coordinate grid + axes: {xmin}{xmax}{ymin}{ymax}
\newcommand{\lgrid}[4]{%
  \draw[step=1, linegray!40, very thin] (#1,#3) grid (#2,#4);
  \draw[->, semithick, charcoal] (#1,0)--(#2,0) node[below right, font=\scriptsize]{$x$};
  \draw[->, semithick, charcoal] (0,#3)--(0,#4) node[left, font=\scriptsize]{$y$};}
% Vocab answer for the key: bold forest term, then the filled definition on a write-line.
% The leading and trailing \par are required: \ansline ends with \dotfill but does NOT end the
% paragraph, so without them each term label gets pulled onto the previous answer's line.
\newcommand{\vocabans}[2]{%
  \par\noindent\textbf{\textcolor{forest}{#1:}}\\[1pt]\ansline{#2}\par}

\begin{document}

\pageheader{Unit 4, Lesson 4.6}{Guided Notes --- Answer Key}
\namedateperiod

\vspace{0.08in}

\begin{objectivebox}
\textbf{By the end of this lesson, I will be able to\dots}
\begin{itemize}\setlength{\itemsep}{2pt}
  \item \textbf{solve} a rational equation by multiplying every term by the LCD and solving the linear or quadratic equation that is left (A2.EI.4b);
  \item \textbf{check} every candidate against the restriction list and \textbf{verify} the survivors --- algebraically and on a graph (A2.EI.4b/c);
  \item \textbf{explain why} a candidate can be extraneous, using the fact that multiplying by an expression that could be zero is not reversible (A2.EI.4d);
  \item \textbf{write} a rational equation to model a situation, solve it, and say what the answer means --- including which answers to reject and why (A2.EI.4a/c).
\end{itemize}
\end{objectivebox}

\vspace{0.05in}

\begin{vocabbox}
\small
Fill in each term as we build it together.
\par\vspace{2pt}   % \par is required: \vocabans's \noindent is a no-op mid-paragraph
\vocabans{Rational equation}{an equation in which the variable appears in a denominator. Its solutions must come from the domain of \emph{every} expression in it, so the excluded values are off limits before the work even starts}
\vocabans{Clearing the denominators}{multiplying \emph{both sides} --- that is, every single term --- by the LCD, so that every denominator cancels and a linear or quadratic equation is left behind}
\vocabans{Candidate vs.\ solution}{a \emph{candidate} is a number the cleared equation produced; it becomes a \emph{solution} only after it survives the check against the restriction list}
\vocabans{Extraneous solution}{a candidate that solves the cleared equation but not the original one. It appears because multiplying by an expression that equals zero is not reversible, so it is always one of the excluded values from Step 1}
\end{vocabbox}

\begin{hookbox}
Two equations. They are the same except for one number in the denominator.
\[ \textbf{(I)}\quad \frac{x^2}{x-2}=\frac{4}{x-2} \qquad\qquad\qquad
   \textbf{(II)}\quad \frac{x^2}{x-3}=\frac{4}{x-3} \]
Multiply both sides of \emph{either} one by its denominator and the very same equation drops out:
$x^2=4$, so the candidates are $x=$ \ans{$2$} and $x=$ \ans{$-2$} --- for \textbf{both}
equations. Now put those two numbers back into the originals and see what actually happens.
\begin{center}
\renewcommand{\arraystretch}{1.6}
\begin{tabularx}{0.96\linewidth}{|C{1.9cm}|C{1.5cm}|C{2.6cm}|C{2.6cm}|X|}
\hline
\rowcolor{forestbg}
\textbf{Equation} & \textbf{Try $x=$} & \textbf{Left side} & \textbf{Right side} &
\textbf{True? False? Neither?} \\ \hline
(I)  & $2$  & \ans{$\frac40$ undef.} & \ans{$\frac40$ undef.} & \ans{neither --- not a statement} \\ \hline
(I)  & $-2$ & \ans{$\frac{4}{-4}=-1$} & \ans{$\frac{4}{-4}=-1$} & \ans{true} \\ \hline
(II) & $2$  & \ans{$\frac{4}{-1}=-4$} & \ans{$\frac{4}{-1}=-4$} & \ans{true} \\ \hline
(II) & $-2$ & \ans{$\frac{4}{-5}$} & \ans{$\frac{4}{-5}$} & \ans{true} \\ \hline
\end{tabularx}
\end{center}
\vspace{-2pt}
Equation (I) has \ans{$1$} solution(s); equation (II) has \ans{$2$}. The algebra could not
tell them apart --- \emph{only the check could.} That is today's whole lesson.
\end{hookbox}

\begin{notesbox}{1. First: Expression or Equation?}
These two lines look almost identical, and what you are allowed to do to them is completely different.

\vspace{4pt}
\begin{center}
\renewcommand{\arraystretch}{1.4}
\begin{tabularx}{\linewidth}{@{}p{0.30\linewidth} X X@{}}
\hline
 & \textbf{Lesson 4.3 --- an \emph{expression}} & \textbf{Lesson 4.6 --- an \emph{equation}} \\
\hline
What it looks like & $\dfrac{6}{x}-\dfrac{2}{x-1}$ & $\dfrac{6}{x}-\dfrac{2}{x-1}=1$ \\
The job & \emph{Combine} it into one fraction & \emph{Solve} it for $x$ \\
What you do with the LCD & Build it: rewrite each piece \emph{over} the LCD & Multiply \emph{every term} by it, so the denominators cancel away \\
What is left at the end & A fraction, with restrictions & Numbers, which must then be \ans{checked} \\
\hline
\end{tabularx}
\end{center}

\vspace{3pt}
\begin{center}
\fbox{\parbox{0.92\linewidth}{\centering\vspace{2pt}
\textbf{You may clear denominators only when there is an equals sign.} An expression has no sides to
multiply. Clearing the denominators of an expression changes its \emph{value}; clearing them in an
equation keeps the equation balanced.\vspace{2pt}}}
\end{center}
\end{notesbox}

\begin{notesbox}{2. The Four Steps --- Always in This Order}
\begin{center}
\renewcommand{\arraystretch}{1.35}
\begin{tabularx}{\linewidth}{@{}c l X@{}}
\hline
\textbf{\#} & \textbf{Step} & \textbf{What you write down} \\
\hline
1 & \textbf{Factor \& restrict} & factor every denominator; list every \ans{excluded} value. \emph{Write the list at the top of your work.} \\
2 & \textbf{Multiply by the LCD} & every term, including any term with no denominator; then cancel \\
3 & \textbf{Solve} & whatever is left --- a \ans{linear} or a \ans{quadratic} equation \\
4 & \textbf{Check} & compare each answer to the Step 1 list; discard any match; verify the survivors \\
\hline
\end{tabularx}
\end{center}

\vspace{3pt}
\begin{center}
\fbox{\parbox{0.92\linewidth}{\centering\vspace{2pt}
\textbf{Clearing the denominators solves a \emph{different} equation. The check is what brings you
back.}\vspace{2pt}}}
\end{center}

\vspace{3pt}
Until Step 4 is done, the numbers you have are only \textbf{candidates}. Step 1 is not optional
bookkeeping --- it is the only thing that makes Step 4 possible, and it must happen \emph{before} you
touch the equation.
\end{notesbox}

\begin{notesbox}{3. Why This Happens --- and Why It Is Not Your Fault}
From the Warm-Up: multiplying both sides of an equation by a \emph{nonzero} number can always be undone
by dividing, so it never changes the answers. Multiplying both sides by \ans{zero} cannot be
undone --- and it turns \emph{every} statement, true or false, into $0=0$.

\vspace{4pt}
When you multiply a rational equation by its LCD, you are not multiplying by a number. You are
multiplying by an \emph{expression}, and at each excluded value that expression is exactly
\ans{zero}. So at those inputs --- and only at those inputs --- the multiplication is not
reversible, and the cleared equation can pick up answers that the original never had.

\vspace{4pt}
\begin{center}
\fbox{\parbox{0.94\linewidth}{\centering\vspace{2pt}
An \textbf{extraneous solution} is a solution of the cleared equation that is \emph{not} a solution of
the original one. \\[3pt]
\textbf{It is never a random number. It is always a value from your Step 1
list.}\vspace{2pt}}}
\end{center}

\vspace{4pt}
That last sentence is worth its own line, because it makes Step 4 quick. You are not re-solving
anything. You are reading two short lists --- \emph{candidates} and \emph{restrictions} --- and looking
for overlap.

\vspace{3pt}
\textbf{And notice what a check on an excluded value actually looks like.} Substituting $x=2$ into
equation (I) of the Hook does not produce a \emph{false} statement. It produces \ans{division by zero
on both sides} --- which is not a statement at all. You cannot check a number into an equation where
the equation does not exist.
\end{notesbox}

\begin{notesbox}{4. Worked Example A --- Nothing Gets Thrown Away}
\[ \frac{6}{x}-\frac{2}{x-1}=1 \]

\textbf{Step 1 --- factor and restrict.} The denominators are $x$ and $x-1$, already factored, so
$x\neq$ \ans{$0$} and $x\neq$ \ans{$1$}. \; LCD $=$ \ans{$x(x-1)$}

\vspace{4pt}
\textbf{Step 2 --- multiply every term by the LCD.} Watch the term on the right: it has no denominator,
and it still gets multiplied.
\[ x(x-1)\cdot\frac{6}{x} \;-\; x(x-1)\cdot\frac{2}{x-1} \;=\; x(x-1)\cdot 1 \]
Cancel: \quad \ans{$6(x-1)$} $-$ \ans{$2x$} $=$ \ans{$x(x-1)$}

\vspace{4pt}
\textbf{Step 3 --- solve what is left.} Distribute, collect everything on one side, and factor.
\[ 6x-6-2x=x^2-x \qquad\Longrightarrow\qquad
   0=x^2-{\color{keyred}\mathbf{5}}\,x+{\color{keyred}\mathbf{6}}
   \qquad\Longrightarrow\qquad
   0=(x-{\color{keyred}\mathbf{2}})(x-{\color{keyred}\mathbf{3}}) \]
Candidates: $x=$ \ans{$2$} and $x=$ \ans{$3$}.

\vspace{4pt}
\textbf{Step 4 --- check.} Restriction list: \ans{$0,\;1$}. \; Is either candidate on it?
\ans{no} \; So \ans{both} candidates survive. Verify both in the \emph{original}:
\begin{center}
\renewcommand{\arraystretch}{1.5}
\begin{tabularx}{0.90\linewidth}{|C{1.6cm}|C{3.4cm}|C{2.0cm}|X|}
\hline
\rowcolor{forestbg}
\textbf{$x$} & \textbf{$\dfrac{6}{x}-\dfrac{2}{x-1}$} & \textbf{Equals $1$?} & \textbf{Verdict} \\ \hline
$2$ & \ans{$3-2=1$} & \ans{yes} & \ans{solution} \\ \hline
$3$ & \ans{$2-1=1$} & \ans{yes} & \ans{solution} \\ \hline
\end{tabularx}
\end{center}
\vspace{-2pt}
Solution set: \ans{$\{2,\,3\}$}
\end{notesbox}

\begin{notesbox}{5. Worked Example B --- One Survives, One Does Not}
\[ \frac{x}{x+1}+\frac{2}{x-1}=\frac{2}{x^2-1} \]

\textbf{Step 1.} Factor the last denominator first: $x^2-1=($\ans{$x-1$}$)($\ans{$x+1$}$)$.
\; Restrictions: $x\neq$ \ans{$1$}, $x\neq$ \ans{$-1$}. \; LCD $=$ \ans{$(x-1)(x+1)$}

\vspace{4pt}
\textbf{Step 2.} Multiply all three terms by the LCD and cancel:
\[ x({\color{keyred}\mathbf{x-1}})+2({\color{keyred}\mathbf{x+1}})={\color{keyred}\mathbf{2}} \]

\vspace{2pt}
\textbf{Step 3.} $x^2-x+2x+2=2 \;\Longrightarrow\; x^2+x=0 \;\Longrightarrow\;
x({\color{keyred}\mathbf{x+1}})=0$. \; Candidates: $x=$ \ans{$0$}, $x=$ \ans{$-1$}.

\vspace{4pt}
\textbf{Step 4.} One candidate is sitting on the restriction list: $x=$ \ans{$-1$} is
\ans{extraneous}. Throw it away --- and notice you did \emph{not} have to substitute anything to know
that. The other candidate, $x=$ \ans{$0$}, is checked in the original: left side
$=\dfrac{0}{1}+\dfrac{2}{-1}=$ \ans{$-2$}, right side $=\dfrac{2}{-1}=$ \ans{$-2$}. \;
Solution set: \ans{$\{0\}$}

\vspace{5pt}
\textbf{Now see it (A2.EI.4b, graphically).} Move everything to one side and combine with 4.3:
\[ \frac{x}{x+1}+\frac{2}{x-1}-\frac{2}{x^2-1}
   =\frac{x(x-1)+2(x+1)-2}{(x-1)(x+1)}=\frac{x^2+x}{(x-1)(x+1)}
   =\frac{x(x+1)}{(x-1)(x+1)}=\frac{x}{x-1},\quad x\neq-1 \]
Solving the equation is the same as asking where that combined function equals \ans{$0$}. Run
Lesson 4.5 on it: the factor $x+1$ cancels, so $x=-1$ is a \ans{hole}; the factor $x-1$ survives,
so $x=1$ is a \ans{vertical asymptote}; and the simplified numerator is zero at $x=$ \ans{$0$}.

\vspace{3pt}
\begin{center}
\begin{minipage}{0.53\linewidth}\centering
\begin{tikzpicture}[scale=0.40]
  \lgrid{-6}{7}{-4}{6}
  \draw[navy, dashed, semithick] (1,-4)--(1,6);
  \draw[navy, dashed, semithick] (-6,1)--(7,1);
  \node[navy, font=\scriptsize, above right] at (1,5.2) {$x=1$};
  \node[navy, font=\scriptsize, below right] at (5.0,1) {$y=1$};
  \begin{scope}
    \clip (-6,-4) rectangle (7,6);
    \draw[very thick, forest, domain=-6:0.8, samples=140, smooth] plot (\x,{1+1/((\x)-1)});
    \draw[very thick, forest, domain=1.25:7, samples=140, smooth] plot (\x,{1+1/((\x)-1)});
  \end{scope}
  \draw[forest, very thick, fill=white] (-1,0.5) circle (4.2pt);
  \node[charcoal, font=\scriptsize, left] at (-1.4,0.9) {hole};
  \fill[charcoal] (0,0) circle (3.4pt);
\end{tikzpicture}
\end{minipage}%
\begin{minipage}{0.45\linewidth}
\small
The solid dot at \ans{$(0,0)$} is where the combined function is zero. That is the \emph{real}
solution.

\vspace{4pt}
The open circle at $\left(-1,\frac12\right)$ is where the extraneous candidate went. The graph has no
point there at all --- so the equation cannot be satisfied there.

\vspace{4pt}
\textbf{Fill in the rule:} a true solution shows up as an \ans{$x$-intercept}; an extraneous candidate
shows up as a \ans{hole} (or a wall).
\end{minipage}
\end{center}
\end{notesbox}

\begin{notesbox}{6. Cross-Multiplying, and the Answer ``No Solution''}
\textbf{The shortcut.} When the equation is a \textbf{proportion} --- one single fraction equal to one
single fraction --- you may cross-multiply. It is not a new rule; it is Step 2 with the LCD
$=$ (bottom left)(bottom right), already cancelled for you. \emph{It does not apply when there are
three terms.}

\vspace{3pt}
\[ \frac{2}{x-1}=\frac{x}{x+2} \]
Restrictions: $x\neq$ \ans{$1$}, $x\neq$ \ans{$-2$}. \; Cross-multiply:
$2({\color{keyred}\mathbf{x+2}})=x({\color{keyred}\mathbf{x-1}})$, so $2x+4=x^2-x$ and
$0=(x-{\color{keyred}\mathbf{4}})(x+{\color{keyred}\mathbf{1}})$. \\[3pt]
Candidates $x=$ \ans{$4$} and $x=$ \ans{$-1$}. On the restriction list? \ans{no} \;
Solution set: \ans{$\{4,\,-1\}$}

\vspace{6pt}
\textbf{And sometimes everything gets thrown away.}
\[ \frac{1}{x-2}+\frac{1}{x+2}=\frac{4}{x^2-4} \]
Restrictions: $x\neq$ \ans{$2,\;-2$}. \; LCD $=$ \ans{$(x-2)(x+2)$}. \; Clearing gives
$(x+2)+(x-2)=4$, so $2x=4$ and the only candidate is $x=$ \ans{$2$}.

\vspace{2pt}
That candidate \emph{is} on the restriction list. Every candidate is extraneous, so the equation has
\ans{no solution}. Write that as your answer --- it is a complete answer, not a failure to finish.
\end{notesbox}

\begin{practicebox}
Solve, showing all four steps. \; $\dfrac{6}{x^2-9}+\dfrac{1}{x+3}=1$
\begin{enumerate}[leftmargin=1.7em, itemsep=6pt]
  \item Factored denominators \ans{$(x-3)(x+3)$ and $x+3$}; \; restrictions $x\neq$ \ans{$3,\;-3$}; \;
        LCD $=$ \ans{$(x-3)(x+3)$}
  \item After multiplying every term by the LCD (do not forget the $1$):
        \; \ans{$6$} $+$ \ans{$(x-3)$} $=$ \ans{$x^2-9$}
  \item Solve: \; $0=x^2-{\color{keyred}\mathbf{1}}\,x-{\color{keyred}\mathbf{12}}
        =(x-{\color{keyred}\mathbf{4}})(x+{\color{keyred}\mathbf{3}})$. \; Candidates \ans{$4,\;-3$}
  \item Check. Extraneous: \ans{$-3$} \; Solution set: \ans{$\{4\}$}
  \item In one sentence, say how you knew the extraneous one was extraneous \emph{without} substituting
        it. \ansline{$-3$ was already on the restriction list from Step 1 --- it makes the original
        denominators zero, so it can never be a solution}
\end{enumerate}
\end{practicebox}

\begin{teachernote}
\textbf{The Hook is the argument, so do not rush it.} The two equations are chosen so that the algebra
is \emph{literally identical} --- both clear to $x^2=4$ --- and the answers differ. That kills, in
advance, the belief students will otherwise carry all period: that a careful enough solver would not
need to check. The row to dwell on is (I) at $x=2$: both sides read $\frac40$. Ask directly, ``is that
row true or false?'' The answer is neither, and that is the point. A number that makes the equation
\emph{undefined} was never a candidate for making it true.

\vspace{3pt}
\textbf{Section 1 is worth its two minutes.} Students who have just spent 4.3 building common
denominators will, under pressure, ``clear'' the denominators of an expression. The line to post is
\emph{you may clear denominators only when there is an equals sign}. Expect to say it again during the
activity.

\vspace{3pt}
\textbf{Section 3 is the A2.EI.4d requirement}, and it is graded language. ``$x=-1$ doesn't work'' is
not a justification. ``$x=-1$ makes a denominator of the original equation zero, so it is not in the
domain; it appeared only because multiplying by $(x-1)(x+1)$ multiplies by zero there, which is not
reversible'' is. Practice the sentence out loud, and reuse it every time an extraneous solution comes
up.

\vspace{3pt}
\textbf{Section 5's graph is where 4.5 pays its debt.} The combined function is
$\frac{x(x+1)}{(x-1)(x+1)}$ --- a cancelling factor and a surviving one --- so yesterday's cancel test
sorts the two candidates for you: the cancelled factor gives the \emph{hole}, and the hole is exactly
the extraneous candidate. Say the pairing plainly: \textbf{true solutions are $x$-intercepts;
extraneous candidates are holes or walls}. Students with graphing technology can verify a solution set
this way in seconds (A2.EI.4c), and a solution that ``disappears'' on the graph is the visual signature
of an extraneous root.

\vspace{3pt}
\textbf{Section 6, two cautions.} First, cross-multiplying is not a separate method --- present it as
Step 2 pre-cancelled, or students will reach for it on three-term equations and lose a term. Second,
\emph{``no solution'' is an answer.} Students will assume they made a mistake and start over. Note for
yourself that there are two different ways to reach it: every candidate is extraneous (this example),
or the cleared equation is itself never true (e.g.\ it reduces to $0=4$). Both are complete answers,
but only the first involves extraneous solutions --- worth distinguishing if a student raises it, and
homework 2(f) puts both on the same page.

\vspace{3pt}
\textbf{Timing.} If you are short, compress Section 4 --- it is the case where nothing goes wrong and
students can finish it from the pattern. Do not compress Section 5.
\end{teachernote}

\end{document}
